\chapter{Gỡ áp lực}
\markboth{Gỡ áp lực}{Ease the Pressure}


\section{Điểm số quan trọng, nhưng nó không phải toàn bộ giá trị của em.}
\begin{truyen}{Điểm số quan trọng, nhưng nó không phải toàn bộ giá trị ...}{Grades are not everything.}
\chuhoa{K}{hi một học sinh bắt đầu nghĩ mình chỉ đáng giá bằng điểm số, nản học sẽ đến rất nhanh. Lời nhắc đúng lúc giúp các em tách kết quả khỏi phẩm giá của chính mình.}
\textit{(When a student begins to think their worth is equal to grades, discouragement comes quickly. The right reminder helps separate performance from personal worth.)}
\end{truyen}

\begin{baihoc}
Kết quả đo một phần năng lực, không đo toàn bộ con người.
\textit{(Results measure part of performance, not the whole person.)}
\end{baihoc}

\begin{ghinhoanh}
\textbf{Short English to remember:}

Grades matter, but they are not you.

\end{ghinhoanh}

\begin{tuhocnhanh}
1. Em có đang đánh đồng bản thân với điểm số không? \textit{(Are you equating yourself with your grades?)}

2. Ngoài điểm số, em còn điều gì đáng quý? \textit{(Besides grades, what else is valuable in you?)}
\end{tuhocnhanh}
\ngancach


\section{Áp lực không phải lúc nào cũng làm người ta giỏi hơn.}
\begin{truyen}{Áp lực không phải lúc nào cũng làm người ta giỏi hơn.}{Pressure has limits.}
\chuhoa{C}{ó mức ép khiến học sinh tập trung hơn, nhưng cũng có mức ép làm các em tê liệt. Thầy cô cần biết khi nào phải đẩy và khi nào phải tháo bớt.}
\textit{(Some pressure helps students focus, but too much pressure freezes them. Teachers need to know when to push and when to release.)}
\end{truyen}

\begin{baihoc}
Áp lực đúng mức giúp trưởng thành; áp lực quá mức làm học sinh gãy.
\textit{(Healthy pressure helps growth; too much pressure breaks students.)}
\end{baihoc}

\begin{ghinhoanh}
\textbf{Short English to remember:}

Too much pressure shuts the mind down.

\end{ghinhoanh}

\begin{tuhocnhanh}
1. Áp lực nào đang làm em nặng nhất? \textit{(Which pressure weighs on you the most?)}

2. Điều gì sẽ giúp nó vừa sức hơn? \textit{(What would make it more manageable?)}
\end{tuhocnhanh}
\ngancach


\section{Em không cần chứng minh giá trị của mình trong một tuần.}
\begin{truyen}{Em không cần chứng minh giá trị của mình trong một tuần.}{Worth takes time.}
\chuhoa{C}{ó học sinh sống như thể mọi thứ phải được chứng minh ngay lập tức. Nhưng sự trưởng thành thật đi theo tháng năm, không theo một tuần căng thẳng.}
\textit{(Some students live as if everything must be proven immediately. But real growth unfolds over months and years, not one stressful week.)}
\end{truyen}

\begin{baihoc}
Đừng bắt mình hoàn thành cả tương lai chỉ trong vài ngày mệt mỏi.
\textit{(Do not force yourself to complete your whole future in a few exhausting days.)}
\end{baihoc}

\begin{ghinhoanh}
\textbf{Short English to remember:}

You do not need to prove everything this week.

\end{ghinhoanh}

\begin{tuhocnhanh}
1. Em có đang tự ép mình phải giỏi quá nhanh không? \textit{(Are you forcing yourself to improve too fast?)}

2. Em nên cho mình thêm thời gian ở việc gì? \textit{(Where should you give yourself more time?)}
\end{tuhocnhanh}
\ngancach


\section{Mệt thì nghỉ cho đúng, đừng nghỉ kiểu bỏ luôn.}
\begin{truyen}{Mệt thì nghỉ cho đúng, đừng nghỉ kiểu bỏ luôn.}{Rest, don't quit.}
\chuhoa{N}{ghỉ ngơi là cần, nhưng nhiều học sinh vì quá mệt mà bỏ hẳn nhịp học. Một lời hướng dẫn đúng giúp các em nghỉ để hồi sức chứ không nghỉ để trôi luôn.}
\textit{(Rest is necessary, but many tired students stop studying altogether. The right guidance helps them rest to recover, not rest until they drift away.)}
\end{truyen}

\begin{baihoc}
Nghỉ đúng cách là một phần của việc đi đường dài.
\textit{(Resting well is part of going the distance.)}
\end{baihoc}

\begin{ghinhoanh}
\textbf{Short English to remember:}

Rest wisely. Return stronger.

\end{ghinhoanh}

\begin{tuhocnhanh}
1. Em nghỉ ngơi kiểu nào thường làm mình đỡ hơn? \textit{(What kind of rest actually helps you recover?)}

2. Khi nào em nghỉ quá lâu thành chán luôn? \textit{(When does your break become avoidance?)}
\end{tuhocnhanh}
\ngancach


\section{Có những lúc bớt một mục tiêu lại là cứu được một học sinh.}
\begin{truyen}{Có những lúc bớt một mục tiêu lại là cứu được một học sinh.}{Less can save more.}
\chuhoa{K}{hông phải lúc nào thêm yêu cầu cũng là tốt. Có khi bớt một việc đúng lúc lại giúp học sinh giữ được niềm tin để đi tiếp phần còn lại.}
\textit{(Adding more demands is not always better. Sometimes removing one thing at the right time helps a student keep enough hope to continue.)}
\end{truyen}

\begin{baihoc}
Giảm bớt đúng lúc không phải nuông chiều; đó là biết ưu tiên điều quan trọng hơn.
\textit{(Reducing the load at the right time is not indulgence; it is wise prioritizing.)}
\end{baihoc}

\begin{ghinhoanh}
\textbf{Short English to remember:}

Less pressure can save more effort.

\end{ghinhoanh}

\begin{tuhocnhanh}
1. Nếu được bớt một áp lực, em muốn bớt điều gì? \textit{(If one pressure could be reduced, which would you choose?)}

2. Điều gì thật sự quan trọng nhất lúc này? \textit{(What truly matters most right now?)}
\end{tuhocnhanh}
\ngancach


\section{Em không cần mang mọi kỳ vọng của người lớn lên vai một mình.}
\begin{truyen}{Em không cần mang mọi kỳ vọng của người lớn lên vai một m...}{Too many expectations crush.}
\chuhoa{C}{ó học sinh học trong nỗi sợ làm người khác thất vọng. Một câu nói giúp các em thở ra rằng: em được phép là con người trước khi là kỳ vọng của ai đó.}
\textit{(Some students study in fear of disappointing others. One sentence can help them breathe again: you are allowed to be a person before becoming someone's expectation.)}
\end{truyen}

\begin{baihoc}
Kỳ vọng cần nâng đỡ, không nên đè lên vai học sinh đến mức nghẹt thở.
\textit{(Expectations should support students, not crush them.)}
\end{baihoc}

\begin{ghinhoanh}
\textbf{Short English to remember:}

You are more than other people's expectations.

\end{ghinhoanh}

\begin{tuhocnhanh}
1. Kỳ vọng nào làm em mệt nhất? \textit{(Which expectation tires you the most?)}

2. Em cần nói gì với bản thân để nhẹ hơn? \textit{(What do you need to tell yourself to feel lighter?)}
\end{tuhocnhanh}
\ngancach


\section{Không sao nếu em phải học lại cách bắt đầu.}
\begin{truyen}{Không sao nếu em phải học lại cách bắt đầu.}{Start over calmly.}
\chuhoa{C}{ó những giai đoạn học sinh mất nhịp hoàn toàn. Lúc đó, điều cần nhất không phải là xấu hổ, mà là học lại cách bắt đầu từ những bước rất nhỏ.}
\textit{(Some periods throw students completely off rhythm. What they need then is not shame, but to learn how to start again from very small steps.)}
\end{truyen}

\begin{baihoc}
Bắt đầu lại không làm em kém; nó chứng tỏ em chưa bỏ mình.
\textit{(Starting again does not make you lesser; it proves you have not abandoned yourself.)}
\end{baihoc}

\begin{ghinhoanh}
\textbf{Short English to remember:}

Starting again is still courage.

\end{ghinhoanh}

\begin{tuhocnhanh}
1. Em đang cần bắt đầu lại điều gì? \textit{(What do you need to restart?)}

2. Bước nhỏ đầu tiên của em là gì? \textit{(What is your first small step?)}
\end{tuhocnhanh}
\ngancach


\section{Thầy không muốn em kiệt sức để đổi lấy một con điểm đẹp.}
\begin{truyen}{Thầy không muốn em kiệt sức để đổi lấy một con điểm đẹp.}{Health before score.}
\chuhoa{M}{ột con điểm đẹp mà đổi bằng sự kiệt sức kéo dài không phải chiến thắng trọn vẹn. Người thầy tốt phải biết nhắc học sinh giữ mình trước khi giữ điểm.}
\textit{(A nice score earned by long exhaustion is not a complete victory. A good teacher reminds students to protect themselves before protecting grades.)}
\end{truyen}

\begin{baihoc}
Học để lớn lên, không phải để mòn đi.
\textit{(Study to grow, not to wear yourself away.)}
\end{baihoc}

\begin{ghinhoanh}
\textbf{Short English to remember:}

Do not trade your health for one score.

\end{ghinhoanh}

\begin{tuhocnhanh}
1. Em có đang mệt quá mức vì học không? \textit{(Are you too exhausted because of study?)}

2. Em cần chỉnh điều gì để bền hơn? \textit{(What should you adjust to last longer?)}
\end{tuhocnhanh}
\ngancach


\section{Một lời dịu đi đúng lúc có thể cứu cả ý chí học tập.}
\begin{truyen}{Một lời dịu đi đúng lúc có thể cứu cả ý chí học tập.}{Gentleness restores energy.}
\chuhoa{K}{hi học sinh đang quá căng, thêm một lời gắt dễ làm mọi thứ gãy luôn. Trái lại, một lời dịu đúng lúc có thể giữ lại phần ý chí còn sót lại.}
\textit{(When students are under too much pressure, one more harsh word can break everything. A gentle word at the right time can save what is left of their will.)}
\end{truyen}

\begin{baihoc}
Sự dịu dàng đúng lúc không làm việc học yếu đi; nó làm lòng người bớt sập.
\textit{(Timely gentleness does not weaken learning; it keeps hearts from collapsing.)}
\end{baihoc}

\begin{ghinhoanh}
\textbf{Short English to remember:}

Gentleness can protect effort.

\end{ghinhoanh}

\begin{tuhocnhanh}
1. Lời dịu dàng nào em từng nhớ rất lâu? \textit{(What gentle sentence have you remembered for a long time?)}

2. Khi em căng thẳng, em cần được nói với mình thế nào? \textit{(How do you need people to speak to you when you are stressed?)}
\end{tuhocnhanh}
\ngancach


\section{Muốn học sinh đi xa, đôi khi phải gỡ bớt đá khỏi balô của các em.}
\begin{truyen}{Muốn học sinh đi xa, đôi khi phải gỡ bớt đá khỏi balô của ...}{Travel lighter.}
\chuhoa{K}{hông phải em nào nản cũng vì thiếu ý chí. Có em chỉ đang mang quá nhiều thứ cùng lúc. Giúp các em đi xa hơn đôi khi bắt đầu bằng việc làm gánh nặng nhẹ đi.}
\textit{(Not every discouraged student lacks willpower. Some are simply carrying too much at once. Helping them go farther may begin by lightening the load.)}
\end{truyen}

\begin{baihoc}
Gỡ bớt gánh nặng là cách giữ lại sức đi đường.
\textit{(Lightening the burden is how you preserve the strength to continue.)}
\end{baihoc}

\begin{ghinhoanh}
\textbf{Short English to remember:}

Lighter hearts go farther.

\end{ghinhoanh}

\begin{tuhocnhanh}
1. `Viên đá' nào đang nặng nhất trong balô của em? \textit{(What is the heaviest `stone' in your backpack right now?)}

2. Em có thể nhờ ai giúp gỡ bớt nó? \textit{(Who can help you lighten it?)}
\end{tuhocnhanh}
\ngancach